% Options for packages loaded elsewhere
\PassOptionsToPackage{unicode}{hyperref}
\PassOptionsToPackage{hyphens}{url}
\PassOptionsToPackage{dvipsnames,svgnames*,x11names*}{xcolor}
%
\documentclass[
  10pt,
]{article}
\usepackage{lmodern}
\usepackage{amssymb,amsmath}
\usepackage{ifxetex,ifluatex}
\ifnum 0\ifxetex 1\fi\ifluatex 1\fi=0 % if pdftex
  \usepackage[T1]{fontenc}
  \usepackage[utf8]{inputenc}
  \usepackage{textcomp} % provide euro and other symbols
\else % if luatex or xetex
  \usepackage{unicode-math}
  \defaultfontfeatures{Scale=MatchLowercase}
  \defaultfontfeatures[\rmfamily]{Ligatures=TeX,Scale=1}
\fi
% Use upquote if available, for straight quotes in verbatim environments
\IfFileExists{upquote.sty}{\usepackage{upquote}}{}
\IfFileExists{microtype.sty}{% use microtype if available
  \usepackage[]{microtype}
  \UseMicrotypeSet[protrusion]{basicmath} % disable protrusion for tt fonts
}{}
\makeatletter
\@ifundefined{KOMAClassName}{% if non-KOMA class
  \IfFileExists{parskip.sty}{%
    \usepackage{parskip}
  }{% else
    \setlength{\parindent}{0pt}
    \setlength{\parskip}{6pt plus 2pt minus 1pt}}
}{% if KOMA class
  \KOMAoptions{parskip=half}}
\makeatother
\usepackage{xcolor}
\IfFileExists{xurl.sty}{\usepackage{xurl}}{} % add URL line breaks if available
\IfFileExists{bookmark.sty}{\usepackage{bookmark}}{
\usepackage{hyperref}
}
\hypersetup{
  pdftitle={TP sur les tableaux},
  pdfauthor={Première NSI, Lycée du Parc},
  colorlinks=true,
  linkcolor=Maroon,
  filecolor=Maroon,
  citecolor=Blue,
  urlcolor=Blue,
  pdfcreator={LaTeX via pandoc}}
\urlstyle{same} % disable monospaced font for URLs
\usepackage[top=20mm,left=20mm,right=20mm,heightrounded]{geometry}
\usepackage{listings}
\newcommand{\passthrough}[1]{#1}
\lstset{defaultdialect=[5.3]Lua}
\lstset{defaultdialect=[x86masm]Assembler}
\setlength{\emergencystretch}{3em} % prevent overfull lines
\providecommand{\tightlist}{%
  \setlength{\itemsep}{0pt}\setlength{\parskip}{0pt}}
\setcounter{secnumdepth}{5}

\title{TP sur les tableaux}
\usepackage{etoolbox}
\makeatletter
\providecommand{\subtitle}[1]{% add subtitle to \maketitle
  \apptocmd{\@title}{\par {\large #1 \par}}{}{}
}
\makeatother
\subtitle{Thème types construits}
\author{Première NSI, \href{https://frederic-junier.org/}{Lycée du
Parc}}
\date{}

%%%jolis boites

\usepackage{fancybox, graphicx}



%%%%%%%%%%%%%%%%Packages et Macros Frederic%%%%%%%%%%%%%%%%%%%%%%%%%%%%%


%%%%Insertion de liens hypertextes %%%%

            
%%%%%%%%%%PSTricks%%%%%%%%%%%%

\usepackage{pstricks,pst-plot,pst-text,pst-tree,pst-eps,pst-fill,pst-node,pst-math,pstricks-add,pst-xkey,pst-eucl}


%%%%%%%Tikz%%%%%%%%%%%%%%%
\usepackage{pgf,tikz,tkz-tab}
% Pour les tableaux de signes ou de variations avec tkz-tab voir https://zestedesavoir.com/tutoriels/439/des-tableaux-de-variations-et-de-signes-avec-latex/#1-13389_tikz-un-package-qui-en-a-dans-le-ventre
\usetikzlibrary{arrows}
\usetikzlibrary{shapes.geometric}
\usetikzlibrary{shapes.geometric}
\usetikzlibrary{petri}
\usetikzlibrary{decorations}
\usetikzlibrary{arrows}
\usetikzlibrary{math}
 %Variables must be declared in a tikzmath environment but
       % can be used outside
%       \tikzmath{int \n; \n = 508; \x1 = 1; \y1 =1; 
%                   %computations are also possible
%                    \x2 = \x1 + 1; \y2 =\y1 +3; } 


%%%%%%%%%%%%%%%%%%%%%%%%%%%%%%%%%%%%%%%%
%%%%%%%%%%%Commandes Tikz Perso%%%%%%%%%%%%%%%

% Définition des nouvelles options xmin, xmax, ymin, ymax
% Valeurs par défaut : -3, 3, -3, 3
\tikzset{
xmin/.store in=\xmin, xmin/.default=-3, xmin=-3,
xmax/.store in=\xmax, xmax/.default=3, xmax=3,
ymin/.store in=\ymin, ymin/.default=-3, ymin=-3,
ymax/.store in=\ymax, ymax/.default=3, ymax=3,
}
% Commande qui trace la grille entre (xmin,ymin) et (xmax,ymax)
\newcommand {\grille}[2]
{\draw[help lines,black, thick] (\xmin,\ymin) grid[xstep=#1, ystep=#2] (\xmax,\ymax);}
% Commande \axes
\newcommand {\axes} {
\draw[->,very thick] (\xmin,0) -- (\xmax,0);
\draw[->,very thick] (0,\ymin) -- (0,\ymax);
\draw (0.95*\xmax, 0) node[above] {};
\draw (0, 0.95*\ymax) node[left] {};
}
% Commande qui limite l?affichage à (xmin,ymin) et (xmax,ymax)
\newcommand {\fenetre}
{\clip (\xmin,\ymin) rectangle (\xmax,\ymax);}

%Exemple d'utilisation

%\begin{center}
%\begin{tikzpicture} [xmin=-2,xmax=2,ymin=0,ymax=5]
%\grille{1} \axes \fenetre
%\draw plot[smooth] (\x,\x^2);
%\end{tikzpicture}
%\end{center}

%style pour la perspective cavalière française
%voir Tikz pour l'impatient page 68
\tikzset{math3d/.style=
{x= {(-0.353cm,-0.353cm)}, z={(0cm,1cm)},y={(1cm,0cm)}}}

%%%%%%%Symbole pour code calculatrice%%%%%%

%Flèche remplie pour défilement de menu

\newcommand{\flechefillright}{
\begin{tikzpicture}[scale=0.15] \fill (0,0)--(2,1)--(0,2)--cycle;
\end{tikzpicture}}

%%%%%%%%%%%%%Symboles pour calculatrice Casio%%%%
\newcommand{\execasio}{\Pisymbol{psy}{191}} %Retour chariot
\newcommand{\dispcasio}{\begin{pspicture}(.1,.1)\pspolygon*(.1,0)(.1,.1)\end{pspicture}} %Triangle « Disp »
\newcommand{\dispcasiotikz}{\begin{tikzpicture}[scale=0.2]
\fill (0,0) -- (1,0) -- (1,1) -- cycle;
\end{tikzpicture}} %Triangle « Disp »
%

%Fleche entre deux lignes, d'apres 'un bon petit' : http://forum.mathematex.net/latex-f6/fleches-entre-deux-lignes-pour-resolution-d-equation-t10283.html#p99817
\newcommand\addnode[1]{\Rnode{#1}{}}
\newcommand\linknode[3]{\ncbar[angleA=0,angleB=0,nodesep=1ex,arm=10ex,offset=-2pt]{->}{#1}{#2}\Aput{\vphantom{x}#3}}


%%Commande pour touche de calculatrice

\newcommand\tc[1]{%
{
\begin{tikzpicture}
\node[draw,rectangle,rounded corners=3pt] (P) at (0,0){#1};
\end{tikzpicture}
}
}

%%%%%%%%%%%%%%%%%%%%%%%%%%%%%%%%%%%%%%%%
%%%%%%%%%%%Fin Commandes Tikz%%%%%%%%%%%%%%%


%%%%%%%%%%%%Specifiques%%%%%%%%%%%
\usepackage{wrapfig}
%pour insérer une figure à droite ou à gauche d'un texte
%\begin{wrapfigure}[nb lignes]{placement l,r,c,i(inside),o(outside)}[overhang]{width}
%ce package fonctionne mal à proximité des listes
%%%%%%%%%%%%%%%%%%%%%%%%%%%%%%%%%%%%%

%%%%%Environnements et symboles spéciaux pour faire joli%%%%%%

%%%Bclogo, pour des environnements + jolis avec insertion de logo%%%%
%Dépendances de  bclogo
\usepackage{xkeyval}  
\usepackage{etoolbox}
\usepackage{ifpdf}
\usepackage[framemethod=tikz]{mdframed}
\usepackage[tikz]{bclogo}

%\newcommand\bcpython{\includegraphics[width=17pt]{/home/fjunier/Maths/python-logo.eps}}
\newcommand\bcpython{\includegraphics[width=17pt]{/home/fjunier/Maths/python-logo.png}}
%\newcommand\bcpython{\includegraphics[width=17pt]{/home/frederic/Maths/python-logo.png}}

%% Framed
\usepackage{framed}  %Le package « framed» Crée 3 nouveaux environnements, qui se comportent comme des minipage de largeur \linewidth, mais permettant en plus de se casser entre plusieurs pages.     * framed : avec un cadre autour;     * shaded : avec un fonc coloré (il faut définir la couleur shadecolor);     * leftbar : avec une barre le long du côté gauche.

%%%%%%%%%%%%%%%%%%%Présentation de codes sources%%%%%%%%%%%%%%%%%
\usepackage{listings}
%On utilise l?environnement lstlisting pour insérer
%un code source.
%En plus de l?environnement lstlisting, on peut également utiliser la
%commande \lstinline qui fonctionne comme la commande \verb, en ce
%sens qu?on peut utiliser n?importe quel caractère comme délimiteur. Enfin,
%la commande \lstinputlisting permet de charger un code source depuis
%un fichier externe.
%Il y a deux manières de préciser des options : soit via l?option de l?envi-
%ronnement ou de la commande, soit en utilisant la commande \lstset
%qui permet de définir des options de manière globale.

\lstset{ %
  language=Python,                % the language of the code
  basicstyle=\ttfamily,           % the size of the fonts that are used for the code
  %numbers=left,                   % where to put the line-numbers
  numberstyle=\tiny,  % the style that is used for the line-numbers
  %stepnumber=2,                   % the step between two line-numbers. If it's 1, each line 
                                  % will be numbered
  %numbersep=5pt,                  % how far the line-numbers are from the code
  backgroundcolor=\color{white},      % choose the background color. You must add \usepackage{color}
  showspaces=false,               % show spaces adding particular underscores
  showstringspaces=false,         % underline spaces within strings
  showtabs=false,                 % show tabs within strings adding particular underscores
  frame=single,                   % adds a frame around the code
  rulecolor=\color{black},        % if not set, the frame-color may be changed on line-breaks within not-black text (e.g. comments (green here))
  tabsize=4,                      % sets default tabsize to 2 spaces
  captionpos=b,                   % sets the caption-position to bottom
  breaklines=true,                % sets automatic line breaking
  breakatwhitespace=false,        % sets if automatic breaks should only happen at whitespace
  %title=\lstname,                   % show the filename of files included with \lstinputlisting;
                                  % also try caption instead of title
  breakindent=1cm,
  keywordstyle=\color{blue},          % keyword style
  commentstyle=\color{red},       % comment style
  %stringstyle=\ttfamily\color{green},         % string literal style
  escapeinside={\%*}{*)},            % if you want to add LaTeX within your code
  morekeywords={*,...},              % if you want to add more keywords to the set
  deletekeywords={...}              % if you want to delete keywords from the given language
  upquote=true,columns=flexible,
xleftmargin=1cm,xrightmargin=1cm,
 inputencoding=utf8,			%Les lignes qui suivent sont pour le codage utf8
  extendedchars=true,
  literate=%
            {é}{{\'{e}}}1
            {è}{{\`{e}}}1
            {ê}{{\^{e}}}1
            {ë}{{\¨{e}}}1
            {û}{{\^{u}}}1
            {ù}{{\`{u}}}1
            {â}{{\^{a}}}1
            {à}{{\`{a} }}1
            {î}{{\^{i}}}1
            {ô}{{\^{o}}}1
            {ç}{{\c{c}}}1
            {Ç}{{\c{C}}}1
            {É}{{\'{E}}}1
            {Ê}{{\^{E}}}1
            {À}{{\`{A}}}1
            {Â}{{\^{A}}}1
            {Î}{{\^{I}}}1
}

\lstdefinestyle{rond}{
  numbers=none,
  backgroundcolor=\color{gristclair},
  frameround =tttt
}

\lstdefinestyle{compil}{
  numbers=none,
  backgroundcolor=\color{gristclair}
}
%\lstset{language=Python,basicstyle=\small , frame=single,tabsize=4,showspaces=false,showtabs=false,showstringspaces=false,numbers=left,numberstyle=\tiny , extendedchars=true}



%%%%%%%%%%%%%%%%%%%%%%%%%%%%%%%%%%%%%%%%%%%%%%%%%%%%%%%%%%%%%%%%%%%%%%%%
%%%%%%%%%%%%%%%%%%%%Environnements persos%%%%%%%%%%%%%%%%%%%%%%%%%%%%%%%%
%Syntaxe :
%\newenvironment{nom}[nombre d'args][defaut]{definitions initiales}{definitions finales}
%definitions intiales sont les commandes appelées par \begin{nom}
%Definitions finales sont les commandes appelées par \end{nom}

%%%%%%%%%%%%%%%%Définitions des environnemts persostheoreme, exemple ..%%%%
%%%% Exercice avec encadré %%%%
\newcounter{exo}
\newenvironment{exercice}[1]
{\par \medskip   \addtocounter{exo}{1} \noindent  
\begin{bclogo}[arrondi =0.1,   noborder = true, logo=\bccrayon, marge=4]{~\textbf{Exercice} \textbf{\theexo} {\itshape #1} }  \par}
{
\end{bclogo}
 \par \bigskip }

%%Axiomes, Theoremes, Propriété, Définition, Methode, Preuve


\newenvironment{axiome}[1]
{\par \medskip   \begin{leftbar} \noindent \underline{\textbf{Axiome}}\hspace{0.5cm}{\itshape #1}   \vspace*{10pt} \par }
{\end{leftbar}  \par \medskip }


\newcounter{thme}
\newenvironment{theoreme}[1]
{\par \medskip  \addtocounter{thme}{1} \noindent  
\begin{bclogo}[arrondi =0.1,  ombre = true, barre=none, logo=\bcbook, marge=4]{~\textbf{Théorème} \textbf{\thethme} {\itshape #1} }   \par}
{
\end{bclogo}
 \par \bigskip}

 \newenvironment{theoremedef}[1]
{\par \medskip   \addtocounter{thme}{1} \noindent  
\begin{bclogo}[arrondi =0.1,  ombre = true, barre=none, logo=\bcbook, marge=4]{~\textbf{Théorème-Définition} \textbf{\thethme} {\itshape #1} }   \par}
{
\end{bclogo}
 \par \bigskip }
 
\newcounter{prop}
\newenvironment{propriete}[1]
{\par \medskip   \addtocounter{prop}{1} \noindent  
\begin{bclogo}[arrondi =0.1,  ombre = true, barre=none, logo=\bcbook, marge=4]{~\textbf{Propriété} \textbf{\theprop} {\itshape #1} }   \par}
{
\end{bclogo}
 \par \bigskip }


\newenvironment{corollaire}[1]
{\par \medskip   \noindent  
\begin{bclogo}[arrondi =0.1,  ombre = true, barre=none, logo=\bcbook, marge=4]{~\textbf{Corollaire} {\itshape #1} } \par }
{
\end{bclogo}
 \par \bigskip }

\newenvironment{demo}[1]
{\par \medskip   \noindent  
\begin{bclogo}[arrondi =0.1,  ombre = true, barre=zigzag, noborder = true, logo=\bcloupe, marge=0]{~\textbf{Démonstration} {\itshape #1} } \par \vspace{10pt}}
{
\end{bclogo}
 \par \bigskip }

\newcounter{activite}
\newenvironment{activite}[1]
{\par \medskip   \noindent   \addtocounter{activite}{1}
\begin{bclogo}[arrondi =0.1,   noborder = true, logo=\bcvelo, marge=4]{~\textbf{Activité} \textbf{\theactivite} {\itshape #1} }  \par}
{
\end{bclogo}
 \par \bigskip }


\newcounter{rque}
\newenvironment{remarque}
{\par \medskip    \addtocounter{rque}{1} \noindent  
\begin{bclogo}[arrondi =0.1,  ombre = true, barre=snake, noborder = true, logo=\bcinfo, marge=0]{~\textbf{Remarque} \textbf{\therque}}  \par }
{
\end{bclogo}
 \par \bigskip }

\newcounter{def}
\newenvironment{definition}[1]
{\par \medskip   \addtocounter{def}{1} \noindent  
\begin{bclogo}[arrondi =0.1,  ombre = true, barre=none, logo=\bcbook, marge=4]{~\textbf{Définition} \textbf{\thedef} {\itshape #1} }  \par}
{
\end{bclogo}
 \par \bigskip }
 
 
 \newcounter{cours}
\newenvironment{cours}[1]
{\par \medskip   \addtocounter{cours}{1} \noindent  
\begin{bclogo}[arrondi =0.1,  ombre = true, barre=none, logo=\bcbook, marge=4]{~\textbf{Point de cours} \textbf{\thecours} {\itshape #1} }  \par}
{
\end{bclogo}
 \par \bigskip }
 
 

\newenvironment{introduction}
{\par \medskip    \noindent  
 \begin {bclogo}[couleur = blue!5 , arrondi =0.1,logo=\bcrosevents, marge=4] {~\textbf{Introduction}    }
 \par }
{
\end{bclogo}
 \par \bigskip }
 
\newenvironment{memo}[1]
{\par \medskip    \noindent  
\begin{bclogo}[arrondi =0.1,  ombre = true, barre=none, logo=\bccle, marge=4]{~\textbf{À retenir}  {\itshape #1} }  \par}
{
\end{bclogo}
 \par \bigskip }
 
\newcounter{exple}
\newenvironment{exemple}[1]
{\par \medskip   \addtocounter{exple}{1} \noindent  
\begin{bclogo}[arrondi =0.1,   noborder = true, logo=\bclampe, marge=4]{~\textbf{Exemple} \textbf{\theexple} {\itshape #1} }  \par}
{
\end{bclogo}
 \par \bigskip }




\newcounter{alg}
\newenvironment{algorithme}[1]
{\par \medskip   \addtocounter{alg}{1} \noindent  
 \begin {bclogo}[noborder = true, barre=zigzag,logo=\bcpython, marge=4] {~\textbf{Algorithmique} \textbf{\thealg} {\itshape #1} }  \par}
{
\end{bclogo}
 \par \bigskip }

\newcounter{prog}
\newenvironment{programme}[1]
{\par \medskip   \addtocounter{prog}{1} \noindent  
 \begin {bclogo}[noborder = true, barre=zigzag,logo=\bcpython, marge=4] {~\textbf{Programme} \textbf{\theprog} {\itshape #1} }  \par  \bigskip}
{
\end{bclogo}
 \par \bigskip }
 
\newcounter{logi}
\newenvironment{logique}[1]
{\par \medskip   \addtocounter{logi}{1} \noindent  
 \begin {bclogo}[noborder = true, barre=zigzag,logo=\bclampe, marge=4] {~\textbf{Logique} \textbf{\thelogi} {\itshape #1} }  \par}
{
\end{bclogo}
 \par \bigskip }


\newenvironment{methode}[1]
{\par \medskip    \noindent  
 \begin {bclogo}[arrondi =0.1,logo=\bcoutil, marge=4,noborder = true] {~\textbf{Méthode}   {\itshape #1} }  \par}
{
\end{bclogo}
 \par \bigskip }


\newcounter{histo}
\newenvironment{histoire}[1]
{\par \medskip   \addtocounter{histo}{1} \noindent  
 \begin {bclogo}[couleur = blue!10 , arrondi =0.1,logo=\bchorloge, marge=4] {~\textbf{Histoire} \textbf{\thehisto} {\itshape #1} }  \par}
{
\end{bclogo}
 \par \bigskip }




%Environnement contenu pour un document présentant une progression annuelle
\newenvironment{contenu}
{\par \medskip   \begin {bclogo}[ noborder = true,logo=\bccrayon] \noindent {\large \textbf{Contenu de la séance}} \vspace*{10pt} \par  }
{\end{bclogo}  \par \medskip }

%Environnement programme pour un document présentant une progression annuelle
%\newenvironment{programme}
%{\par \medskip   \begin {bclogo}[ noborder = true, barre=zigzag,logo=\bcinfo] \noindent {\large \textbf{Programme officiel}} \vspace*{10pt} \par  }
%{\end{bclogo}  \par \medskip }

%Environnement programme pour un document présentant une progression annuelle
\newenvironment{ressource}
{\par \medskip   \begin {bclogo}[ noborder = true,logo=\bcbook] \noindent {\large \textbf{Ressources}}\\vspace*{10pt} \par }
{\end{bclogo}  \par \medskip }




%%%%%%%%%%%%%%%%%%Maths divers%%%%%%%%%%%%%%%%%%%%%%%%%
%%%%%%%%%%%%%Nombres%%%%%%%%%%%%%%%%

%Ensemble prive de...
%\newcommand{\prive}{\boi}%{\backslash}

%Ensembles de nombres%%%%%%%%%%%%%%%%%
\newcommand{\R}{\mathbb{R}}
\newcommand{\N}{\mathbb{N}}
\newcommand{\D}{\mathbb{D}}
\newcommand{\Z}{\mathbb{Z}}
\newcommand{\Q}{\mathbb{Q}}
%\newcommand{\C}{\mathbb{C}}
\newcommand{\df}{~\ensuremath{]0;+\infty[}~}
\newcommand{\K}{\mathbb{K}}

%%%%%%%%Arithmetique%%%%%%%%%%
%PGCD, PPCM
\newcommand{\PGCD}{\mathop{\rm PGCD}\nolimits}
\newcommand{\PPCM}{\mathop{\rm PPCM}\nolimits}

%Intervalles
\newcommand{\interoo}[2]{]#1\, ;\, #2[}
\newcommand{\Interoo}[2]{\left]#1\, ;\, #2\right[}
\newcommand{\interof}[2]{]#1\, ;\, #2]}
\newcommand{\Interof}[2]{\left]#1\, ;\, #2\right]}
\newcommand{\interfo}[2]{[#1\, ;\, #2[}
\newcommand{\Interfo}[2]{\left[#1\, ;\, #2\right[}
\newcommand{\interff}[2]{[#1\, ;\, #2]}
\newcommand{\Interff}[2]{\left[#1\, ;\, #2\right]}
%\newcommand\interentiers #1#2{[\! [#1\, ;\, #2]\! ]}
\newcommand{\interentiers}[2]{\llbracket #1\, ;\, #2\rrbracket}
%


%%%%%%%%%%%%%%Nombres complexes%%%%%

\newcommand{\ic}{\text{i}}
%\newcommand{\I}{\text{i}}
\newcommand{\im}[1]{\text{Im}\left(#1\right)}
\newcommand{\re}[1]{\text{Re}\left(#1\right)}
\newcommand{\Arg}[1]{\text{arg}\left(#1\right)}
\newcommand{\Mod}[1]{\left[#1\right]}
%Parties entière, réelle, imaginaire, nombre i
\newcommand{\ent}[1]{\text{E}\left(#1\right)}
\renewcommand{\Re}{\mathop{\rm Re}\nolimits}
\renewcommand{\Im}{\mathop{\rm Im}\nolimits}
\renewcommand{\i}{\textrm{i}}

%%%%%%%%%%%Probabilites et statistiques%%%%%
\newcommand{\loibinom}[2]{\mathcal{B}\left(#1\ ; \ #2 \right)}
\newcommand{\loinorm}[2]{\mathcal{N}\left(#1\ ; \ #2 \right)}
\newcommand{\loiexp}[1]{\mathcal{E}\left(#1\right)}
\newcommand{\proba}[1]{\mathbb{P}\big(#1\big)}
\newcommand{\probacond}[2]{\mathbb{P}_{#2}\big(#1\big)}
\newcommand{\esperance}[1]{\mathbb{E}\left(#1\right)}
\newcommand{\variance}[1]{\mathbb{V}\left(#1\right)}
\newcommand{\ecart}[1]{\sigma\left(#1\right)}
\newcommand{\dnormx}{\frac{1}{\sqrt{2\pi}} \text{e}^{-\frac{x^2}{2}}}
\newcommand{\dnormt}{\frac{1}{\sqrt{2\pi}} \text{e}^{-\frac{t^2}{2}}}

%Covariance
\newcommand{\cov}{\mathop{\rm cov}\nolimits}
%


%%%%%%%%%%Analyse%%%%%%%%%%%

%%%%%%%%%%%Courbe%%%%%%%%%%%%
\newcommand{\courbe}[1]{\ensuremath{\mathcal{C}_{#1}}}

%%%%%%%Fonction exponentielle%%%%%
\newcommand{\fe}{~fonction exponentielle~}
\newcommand{\e}{\text{e}}

%Fonction cotangente
\newcommand{\cotan}{\mathop{\rm cotan}\nolimits}
%%%%%%%%%%%%%%%%%%%%%%%%%%%%%%%%%%%%%%%%%
%
%Fonctions hyperboliques
\newcommand{\ch}{\mathop{\rm ch}\nolimits}
\newcommand{\sh}{\mathop{\rm sh}\nolimits}


%%%%%%%%%%%%%%Limites%%%%%%
\newcommand{\limite}[2]{\lim\limits
_{x \to #1} #2}
\newcommand{\limitesuite}[1]{\lim\limits
_{n \to +\infty} #1}
\newcommand{\limiteg}[2]{\lim\limits
_{\substack{x \to #1 \\ x < #1 }} #2}
\newcommand{\limited}[2]{\lim\limits
_{\substack{x \to #1 \\ x > #1 }} #2}

%%%%%%%%%%Continuité%%%%%%%%%%%
\newcommand{\TVI}{théorème des valeurs intermédiaires}

%%%%%%%%%%%Suites%%%%%%%%%%%%
\newcommand{\suite}[1]{\ensuremath{\left(#1_{n}\right)}}
\newcommand{\Suite}[2]{\ensuremath{\left(#1\right)_{#2}}}
%

%%%%%%%%%%%%%%%Calcul intégral%%%%%%
\newcommand{\dx}{\ensuremath{\text{d}x}}		% dx
\newcommand{\dt}{\ensuremath{\text{d}t}}		% dt
\newcommand{\dtheta}{\ensuremath{\text{d}\theta}}		% dtheta
\newcommand{\dy}{\ensuremath{\text{d}y}}		% dy
\newcommand{\dq}{\ensuremath{\text{d}q}}		% dq

%%%Intégrale%%%
\newcommand{\integralex}[3]{\int_{#1}^{#2} #3 \ \dx}
\newcommand{\integralet}[3]{\int_{#1}^{#2} #3 \ \dt}
\newcommand{\integraletheta}[3]{\int_{#1}^{#2} #3 \ \dtheta}

%%%%%Equivalent%%
\newcommand{\equivalent}[1]{\build\sim_{#1}^{}}

%o et O%%%%
\renewcommand{\o}[2]{\build o_{#1\to #2}^{}}
\renewcommand{\O}[2]{\build O_{#1\to #2}^{}}



%%%%%%%%%%%%%%%Geometrie%%%%%%%%%%%%%%%%%%%%%%%

%%%%%%%%%%%%%%%Reperes%%%%%%%%%%%%%%
\def\Oij{\ensuremath{\left(\text{O},~\vect{\imath},~\vect{\jmath}\right)}}
\def\Oijk{\ensuremath{\left(\text{O},~\vect{\imath},~ \vect{\jmath},~ \vect{k}\right)}}
\def\Ouv{\ensuremath{\left(\text{O},~\vect{u},~\vect{v}\right)}}
\renewcommand{\ij}{(\vec\imath\, ;\vec\jmath\,)}
\newcommand{\ijk}{(\vec\imath\, ;\vec\jmath\, ;\vec k\,)}
\newcommand{\OIJ}{(O\,;\, I\,;\, J\,)}
\newcommand{\repere}[3]{\big(#1\, ;\,\vect{#2} ;\vect{#3}\big)}
\newcommand{\reperesp}[4]{\big(#1\, ;\,\vect{#2} ;\vect{#3} ;\vect{#4}\big)}

%%%%%%%%%Coordonnees%%%%%%%%%%%%%%
\newcommand{\coord}[2]{(#1\, ;\, #2)}
\newcommand{\bigcoord}[2]{\big(#1\, ;\, #2\big)}
\newcommand{\Coord}[2]{\left(#1\, ;\, #2\right)}
\newcommand{\coordesp}[3]{(#1\, ;\, #2\, ;\, #3)}
\newcommand{\bigcoordesp}[3]{\big(#1\, ;\, #2\, ;\, #3\big)}
\newcommand{\Coordesp}[3]{\left(#1\, ;\, #2\, ;\, #3\right)}
\newcommand{\Vcoord}[3]{\begin{pmatrix} #1 \\ #2 \\ #3 \end{pmatrix}}
%Symboles entre droites
%\newcommand{\paral}{\sslash}
\newcommand{\paral}{\mathop{/\!\! /}}
%

%%%%%%%%%Produit scalaire, Angles%%%%%%%%%%
\newcommand{\scal}[2]{\vect{#1} \, \cdot \, \vect{#2}}
\newcommand{\Angle}[2]{\left(\vect{#1} \, , \, \vect{#2}\right)}
\newcommand{\Anglegeo}[2]{\left(\widehat{\vect{#1} \, ; \, \vect{#2}}\right)}
\renewcommand{\angle}[1]{\widehat{#1}}
\newcommand{\anglevec}[2]{\left(\vect {#1}\, ,\,\vect {#2} \right)}
\newcommand{\anglevecteur}[2]{(#1\, , \, #2)}
\newcommand{\Anglevec}[2]{(\vecteur{#1}\, ,\,\vecteur{#2})}
\newcommand{\prodscal}[2]{#1 \, \cdot \, #2}
%


%Arc
%\newcommand{\arc}[1]{\wideparen{#1}}
\newcommand{\arcoriente}[1]{\overset{\curvearrowright}{#1}}
%
%


%%%%%%%%%%%%%%%Normes%%%%%%%%%%%%%%%%
\newcommand{\norme}[1]{\left\| #1\right\|}
\newcommand{\normebis}[1]{\delim{2pt}{\|}{9pt}\! #1\delim{2pt}{\|}{9pt}}
\newcommand{\normetriple}[1]{\left |\kern -.07em\left\| #1\right |\kern -.07em\right\|}
\newcommand{\valabs}[1]{\big| \, #1 \, \big|}
%

%%%%%%%%%%%%%%%%%%%%%%%%%%%Degré%%%%%%
%\newcommand{\Degre}{\ensuremath{^\circ}}
%La commande \degre est déjà définie dans le package babel

%%%%%%%%%%Vecteurs%%%%%%%%%%%
\newcommand{\vect}[1]{\mathchoice%
{\overrightarrow{\displaystyle\mathstrut#1\,\,}}%
{\overrightarrow{\textstyle\mathstrut#1\,\,}}%
{\overrightarrow{\scriptstyle\mathstrut#1\,\,}}%
{\overrightarrow{\scriptscriptstyle\mathstrut#1\,\,}}}



%%%%%%%%%%%%%Algebre%%%%%%%%%%%%%%%


%%%%%%%%%%Systemes%%%%%%%%%%%
%Systemes
\newcommand{\sys}[2]{
\left\lbrace
 \begin{array}{l}
  \negthickspace\negthickspace #1\\
  \negthickspace\negthickspace #2\\
 \end{array}
\right.\negthickspace\negthickspace}
\newcommand{\Sys}[3]{
\left\lbrace
 \begin{array}{l}
  #1\\
  #2\\
  #3\\
 \end{array}
\right.}
\newcommand{\Sysq}[4]{
\left\lbrace
 \begin{array}{l}
  #1\\
  #2\\
  #3\\
  #4\\
 \end{array}
\right.}
%
%

%%%%%%%%%%%%%%%%Matrices%%%%%%%%%%%%%%%%%%
%Comatrice
\newcommand{\com}{\mathop{\rm com}\nolimits}
%
%
%Trace
\newcommand{\tr}{\mathop{\rm tr}\nolimits}
%
%
%Transposee
\newcommand{\transposee}[1]{{\vphantom{#1}}^t\negmedspace #1}
%
%
%Noyau
\newcommand{\Ker}{\mathop{\rm Ker}\nolimits}
%
%

%
%Matrices
\newcommand{\Mn}{\mathcal M_n}
\newcommand{\matrice}[4]{
\left(
 \begin{array}{cc}
  #1 & #2 \\
  #3 & #4
 \end{array}
\right)}

\newcommand{\Matrice}[9]{
\left(
 \begin{array}{ccc}
  #1 & #2 & #3\\
  #4 & #5 & #6\\
  #7 & #8 & #9
 \end{array}
\right)}
\newcommand{\Vect}[3]{
\left(\negmedspace
 \begin{array}{c}
  #1\\
  #2\\
  #3
 \end{array}\negmedspace
\right)}
\newcommand{\Ideux}{\matrice{1}{0}{0}{1}}
\newcommand{\Itrois}{\Matrice{1}{0}{0}{0}{1}{0}{0}{0}{1}}
%
%
%Determinants
\newcommand{\determinant}[4]{
\left|
 \begin{array}{cc}
  #1 & #2 \\
  #3 & #4
 \end{array}
\right|}
\newcommand{\Determinant}[9]{
\left|
 \begin{array}{ccc}
  #1 & #2 & #3\\
  #4 & #5 & #6\\
  #7 & #8 & #9
 \end{array}
\right|}

\begin{document}
\maketitle

\emph{On désigne par tableaux, les objets de type
\passthrough{\lstinline!list!} du langage
\href{https://docs.python.org/3/tutorial/datastructures.html}{Python}.
Dans les exercices, comme dans les QCM d'E3C, on rencontrera parfois
l'appellation Python liste pour désigner la structure de données
tableau}.

\begin{exercice}{}

\begin{enumerate}
\def\labelenumi{\arabic{enumi}.}
\tightlist
\item
  Télécharger le matériel de TP depuis
  \url{https://parc-nsi.github.io/premiere-nsi/chapitre6/TP/eleves.zip}
\item
  Extraire l'archive dans un dossier pertinent de son répertoire
  personnel.
\item
  Ouvrir le fichier \passthrough{\lstinline!TP-Tableaux-Eleve.py!} avec
  \passthrough{\lstinline!Spyder!}. Le fichier est prérempli avec des
  cellules contenant les codes à compléter suivis de tests unitaires
  pour tester les fonctions que vous aurez écrites.
\item
  Chaque cellule est préfixée par le commentaire spécial
  \passthrough{\lstinline!\#\%\%!} et peut être exécutée avec
  \passthrough{\lstinline!CTRL + Return!}.
\item
  Écrivez d'abord la fonction avant de la tester avec les tests
  unitaires s'ils sont disponibles : il faut les sélectionner puis les
  décommenter avec \passthrough{\lstinline!CTRL + 1!}.
\end{enumerate}

\end{exercice}

\begin{exercice}{}

QCM type E3C

\begin{enumerate}
\def\labelenumi{\arabic{enumi}.}
\item
  On exécute le code suivant :

\begin{lstlisting}[language=Python]
t = [1,2,3,4,5,6,7,8,9]

v = [c for c in t if c%3 == 0]
\end{lstlisting}
\end{enumerate}

Quelle est la valeur de la variable \passthrough{\lstinline!v!} à la fin
de cette exécution ?

\begin{itemize}
\tightlist
\item
  \textbf{Réponse A :} 18 \textbf{Réponse B :}
  \passthrough{\lstinline![1,4,7]!}
\item
  \textbf{Réponse C :} \passthrough{\lstinline![3,6,9]!} \textbf{Réponse
  D :} \passthrough{\lstinline![1,2,3,4,5,6,7,8,9]!}
\end{itemize}

\begin{enumerate}
\def\labelenumi{\arabic{enumi}.}
\setcounter{enumi}{1}
\tightlist
\item
  On définit :
  \passthrough{\lstinline!resultat = [ i*2 for i in range(10) ]!}.
\end{enumerate}

Quelle est la valeur de \passthrough{\lstinline!resultat!} ?

\begin{itemize}
\tightlist
\item
  \textbf{Réponse A :}
  \passthrough{\lstinline![0, 1, 2, 3, 4, 5, 6, 7, 8, 9]!}
\item
  \textbf{Réponse B :}
  \passthrough{\lstinline![0, 2, 4, 6, 8, 10, 12, 14, 16, 18]!}
\item
  \textbf{Réponse C :}
  \passthrough{\lstinline![0, 2, 4, 6, 8, 10, 12, 14, 16, 18, 20]!}
\item
  \textbf{Réponse D :}
  \passthrough{\lstinline![2, 4, 6, 8, 10, 12, 14, 16, 18]!}
\end{itemize}

\begin{enumerate}
\def\labelenumi{\arabic{enumi}.}
\setcounter{enumi}{2}
\tightlist
\item
  On considère la fonction suivante :
\end{enumerate}

\begin{lstlisting}[language=Python]
def somme(tab):
    s = 0
    for i in range(len(tab)):
        .....
    return s
\end{lstlisting}

Par quelle instruction faut-il remplacer les points de suspension pour
que l'appel \passthrough{\lstinline!somme([10,11,12,13,14])!} renvoie 60
?

\begin{itemize}
\tightlist
\item
  \textbf{Réponse A :} \passthrough{\lstinline!s = tab[i]!}
  \textbf{Réponse B :} \passthrough{\lstinline!s = s + tab[i]!}
\item
  \textbf{Réponse C :} \passthrough{\lstinline!tab[i] = tab[i] + s!}
  \textbf{Réponse D :} \passthrough{\lstinline!s = s + i!}
\end{itemize}

\end{exercice}

\begin{exercice}{}

\begin{enumerate}
\def\labelenumi{\arabic{enumi}.}
\item
  La fonction suivante doit calculer la moyenne d'un tableau de nombres,
  passé en paramètre. Avec quelles expressions faut-il remplacer les
  points de suspension pour que la fonction soit correcte ? \emph{Auteur
  Sylvie Genre}

\begin{lstlisting}[language=Python]
def moyenne(tableau):
    total = ...
    for valeur in tableau:
        total = total + valeur
    return total / ...
\end{lstlisting}
\item
  Donner la valeur des expressions Python suivantes :

\begin{lstlisting}[language=Python]
>>> [1, 2, 3] + [4, 5, 6]
>>> 2 * [1, 2, 3]
\end{lstlisting}
\end{enumerate}

\end{exercice}

\begin{methode}{}

À partir de cet exercice, on précisera pour chaque fonction sa
\textbf{spécification}, dans une chaîne de documentation (ou
\emph{docstring}) multiligne placée juste après l'en-tête de la fonction
entre triples quotes ou guillemets.

Celle-ci est constituée :

\begin{itemize}
\tightlist
\item
  du ou des \textbf{type(s)} du ou des paramètre(s), en précisant si la
  fonction ne prend pas de paramètre
\item
  de la ou des \textbf{précondition(s)} que doivent vérifier ces
  paramètres
\item
  du ou des \textbf{type(s)} de la (ou les) valeur(s) renvoyée(s), en
  précisant si la fonction ne renvoie rien
  (\passthrough{\lstinline!None!} par défaut en Python)
\item
  de la ou des \textbf{postcondition(s)} que doivent vérifier la ou les
  valeur(s) renvoyée(s)
\end{itemize}

Pour vérifier les préconditions on utilise l'instruction
\passthrough{\lstinline!assert!} qui permet de vérifier l'assertion (à
valeur booléenne) qu'elle préfixe : si l'assertion est vraie,
l'exécution continue sinon une exception
\passthrough{\lstinline!AssertionError!} est levée et l'exécution
s'interrompt. On peut passer un message d'erreur personnalisée en le
séparant de l'assertion par une virgule :

\begin{lstlisting}[language=Python]
>>> assert (1 + 1) % 2 == 0, 'attention on calcule modulo 2'
>>> assert (1 + 1) % 2 == 2, 'attention on calcule modulo 2'
Traceback (most recent call last):
  File "<stdin>", line 1, in <module>
AssertionError: attention on calcule modulo 2
\end{lstlisting}

Voici un exemple de rédaction pour la deuxième fonction de l'exercice 4
:

\begin{lstlisting}[language=Python]
def vsom(tab1, tab2):
    """
    Paramètres :
        tab1 et tab2 des tableaux de type list contenant des  nombres de type int ou float
    Préconditions :
        tab1 non vide et len(tab1) == len(tab2)
    Valeur renvoyée :
        un tableau de type list contenant des  nombres de type int ou float
    Postcondition :
        le tableau renvoyé est constitué des sommes
        terme à terme des éléments de tab1 et tab2        
    """
    assert len(tab1) > 0 and len(tab1) == len(tab2); "tab1 et tab2 doivent être de même longueur"
    return ................. #à compléter
\end{lstlisting}

\end{methode}

\begin{exercice}{}

\begin{enumerate}
\def\labelenumi{\arabic{enumi}.}
\tightlist
\item
  Écrire une fonction Python \passthrough{\lstinline!smul!} à deux
  paramètres, un nombre et une liste de nombres, qui multiplie chaque
  élément de la liste par le nombre et renvoie une nouvelle liste :
\end{enumerate}

\begin{lstlisting}[language=Python]
>>> smul(2, [1, 2, 3])
[2, 4, 6]
\end{lstlisting}

\begin{enumerate}
\def\labelenumi{\arabic{enumi}.}
\setcounter{enumi}{1}
\tightlist
\item
  Écrire une fonction Python \passthrough{\lstinline!vsom!} qui prend en
  paramètre deux listes de nombres de même longueur et qui renvoie une
  nouvelle liste constituée de la somme terme à terme de ces deux listes
  :
\end{enumerate}

\begin{lstlisting}[language=Python]
>>> vsom([1, 2, 3], [4, 5, 6])
[5, 7, 9]
\end{lstlisting}

\begin{enumerate}
\def\labelenumi{\arabic{enumi}.}
\setcounter{enumi}{2}
\tightlist
\item
  Écrire une fonction
  \href{https://docs.python.org/3/tutorial/datastructures.html}{Python}
  \passthrough{\lstinline!vdif!} qui prend en paramètre deux listes de
  nombres de même longueur et qui renvoie une nouvelle liste constituée
  de la différence terme à terme de ces deux listes (la deuxième moins
  la première).
\end{enumerate}

\begin{lstlisting}[language=Python]
>>> vdif([1, 2, 3], [4, 5, 6])
[3, 3, 3]
\end{lstlisting}

\begin{enumerate}
\def\labelenumi{\arabic{enumi}.}
\setcounter{enumi}{3}
\tightlist
\item
  Écrire une fonction Python \passthrough{\lstinline!vprod!} qui prend
  en paramètre deux listes de nombres de même longueur et qui retourne
  une nouvelle liste constituée des produits terme à terme de ces deux
  listes.
\end{enumerate}

\begin{lstlisting}[language=Python]
>>> vprod([1, 2, 3], [4, 5, 6]) 
[4, 10, 18]
\end{lstlisting}

\end{exercice}

\begin{exercice}{}

\begin{enumerate}
\def\labelenumi{\arabic{enumi}.}
\tightlist
\item
  Écrire une fonction \passthrough{\lstinline!produit(tab)!} qui
  retourne le produit des éléments d'un tableau de nombres
  \passthrough{\lstinline!tab!}.
\item
  Écrire une fonction \passthrough{\lstinline!all\_positive(tab)!} qui
  retourne un booléen indiquant si tous les éléments du tableau de
  nombres \passthrough{\lstinline!tab!} sont strictement positifs.
\item
  Écrire une fonction \passthrough{\lstinline!any\_positive(tab)!} qui
  retourne un booléen indiquant si au moins un élément du tableau de
  nombres \passthrough{\lstinline!tab!} est strictement positif.
\end{enumerate}

\end{exercice}

\begin{exercice}{}

\begin{enumerate}
\def\labelenumi{\arabic{enumi}.}
\tightlist
\item
  Écrire une fonction
  \passthrough{\lstinline!tableau\_aleatoire(n, a, b)!} qui renvoie un
  tableau de \passthrough{\lstinline!n!} entiers tirés aléatoirement
  entre les entiers \passthrough{\lstinline!a!} et
  \passthrough{\lstinline!b!} inclus. On utilisera la fonction
  \passthrough{\lstinline!randint!} du module
  \passthrough{\lstinline!random!}.
\end{enumerate}

\begin{lstlisting}[language=Python]
Help on method randint in module random:

randint(a, b) method of random.Random instance
    Return random integer in range [a, b], including both end points.
\end{lstlisting}

\begin{enumerate}
\def\labelenumi{\arabic{enumi}.}
\setcounter{enumi}{1}
\item
  Écrire une fonction
  \passthrough{\lstinline!histo\_echantillon(nbexp)!} qui renvoie un
  tableau de taille 6 comptant le nombre d'occurences de chaque face
  numérotée de 1 à 6 sur un échantillon de
  \passthrough{\lstinline!nbexp!} lancers d'un dé cubique équilibré.
\item
  On donne ci-dessous une fonction qui prend en paramètre un tableau et
  renvoie un élément extrait au hasard du tableau. Elle permet par
  exemple de simuler un tirage sans remise d'une boule dans une urne.

\begin{lstlisting}[language=Python]
from random import randint
def tirage_sans_remise(urne):
    """ 
    Paramètres :
        urne un tableau homogène type list
    Précondition :
        urne non vide
    Valeur renvoyée :
        un élément du même type que ceux dans urne
    Postcondition :
        l'élement renvoyée a été extrait aléatoirement de urne
        urne a été modifiée
    """
    return urne.pop(randint(0, .......))
\end{lstlisting}

  \begin{itemize}
  \tightlist
  \item
    Compléter les pointillés pour respecter la spécification de la
    fonction.
  \item
    Écrire une fonction
    \passthrough{\lstinline!echantillon\_tirage\_sans\_remise!}
    respectant la spécification ci-dessous. Les
    \href{https://docs.python.org/fr/3/tutorial/controlflow.html\#function-annotations}{annotations
    de type} sont des métadonnées optionnelles décrivant les types des
    paramètres et des valeurs renvoyées.
  \end{itemize}

\begin{lstlisting}[language=Python]
def echantillon_tirage_sans_remise(urne:list, nbtirage:int) -> list:
    """ 
    Paramètres :
        urne un tableau homogène type list
        nbtirage un entier de type int
    Préconditions :
        nbtirage >= 0
        urne non vide
    Valeur renvoyée :
        un tableau d'entiers de type list
    Postcondition :
        le tableau renvoyé est un échantillon extrait
        aléatoirement de l'urne sans remise
    """
\end{lstlisting}
\end{enumerate}

\end{exercice}

\begin{exercice}{}

\begin{itemize}
\tightlist
\item
  La fonction \passthrough{\lstinline!ord!} prend en paramètre un
  caractère de type \passthrough{\lstinline!string!} et renvoie son
  point de code dans le jeu de caractères
  \href{https://fr.wikipedia.org/wiki/Unicode}{Unicode}.
\item
  La fonction \passthrough{\lstinline!chr!} prend en paramètre un point
  de code Unicode et renvoie le caractère correspondant.
\end{itemize}

\begin{lstlisting}[language=Python]
>>> ord('a')
97
>>> chr(97)
'a'
\end{lstlisting}

\begin{enumerate}
\def\labelenumi{\arabic{enumi}.}
\tightlist
\item
  Construire un tableau \passthrough{\lstinline!alphabet!} qui contient
  toutes les lettres minuscules de l'alphabet romain.
\item
  Construire un tableau \passthrough{\lstinline!consonne!} qui contient
  toutes les consonnes dans \passthrough{\lstinline!alphabet!}.
\item
  Écrire une fonction \passthrough{\lstinline!occurences(chaine)!} qui
  prend en paramètre une chaîne de caractère et renvoie un tableau de
  taille 26 avec le nombre d'occurences dans
  \passthrough{\lstinline!chaine!} des 26 lettres de l'alphabet romain.
\end{enumerate}

\end{exercice}

\begin{exercice}{}

\emph{Exercice du manuel de NSI de Thibault Balabonski chez Ellipses}

En mathématiques, la
\href{http://images.math.cnrs.fr/Mysteres-arithmetiques-de-la-suite-de-Fibonacci.html}{suite
de Fibonacci} est une séquence d'entiers définie ainsi : on part des
entiers 0 et 1 puis on construit à chaque fois l'entier suivant comme la
somme des deux précédents : \[0, 1, 1, 2, 3, 5 ....\]

Compléter la fonction \passthrough{\lstinline!fibonacci(n)!} ci-dessous
pour qu'elle renvoie un tableau contenant les
\passthrough{\lstinline!n!} premiers termes de la suite de Fibonacci
avec \passthrough{\lstinline!n!} entier supposé supérieur ou égal à 2.
Les deux assertions proposées en tests unitaires, doivent être
vérifiées.

\begin{lstlisting}[language=Python]
def fibonacci(n):
    #à compléter

# Tests unitaires (vérification de postconditions)
f6 = fibonacci(6)
assert f6 == [0,1, 1, 2,3,5]
f30 = fibonacci(30)
assert f30[29] == 514229
\end{lstlisting}

\end{exercice}

\begin{methode}{}

Pour stocker de façon persistante (sur le disque et pas seulement en
mémoire vive) des informations lisibles par un humain, on utilise un
\textbf{fichier texte} qui est une simple séquence de caractères. Le
contenu du fichier doit respecter un \textbf{format} qui fixe des
contraintes pour que les informations soient interprétables lors de la
lecture. Par exemple le code source d'un programme
\href{https://docs.python.org/3/tutorial/datastructures.html}{Python} ou
d'une page Web en HTML, un fichier de mesures au format
\href{https://fr.wikipedia.org/wiki/Comma-separated_values}{CSV}, sont
des fichiers textes.

Tester les séquences d'instructions suivantes dans une console
\href{https://docs.python.org/3/tutorial/datastructures.html}{Python}
pour découvrir des manipulations de base de fichiers textes.

\begin{itemize}
\tightlist
\item
  Création d'un fichier \passthrough{\lstinline!test.txt!} dans le
  répertoire de travail
\end{itemize}

\begin{lstlisting}[language=Python]
>>> f = open('test.txt', 'w')   # ouverture du fichier en écriture
>>> f.write('10\n')   #  ligne 1 : '\n'  est un caractère de saut de ligne
3
>>> f.write('11 12\n')  #ligne 2  avec 6 caractères 
6
>>> f.write('13 14 15\n')  #ligne 3 avec 9 caractères
9
>>> f.close()   #fermeture du fichier
\end{lstlisting}

\begin{itemize}
\tightlist
\item
  Lecture et traitement ligne par ligne du fichier
  \passthrough{\lstinline!test.txt!} :
\end{itemize}

\begin{lstlisting}[language=Python]
>>> g = open('test.txt', 'r')  #ouverture du fichier
>>> ligne1 = g.readline()  #lecture de ligne 1
>>> ligne1
'10\n'
>>> ligne1.rstrip() #nettoyage du caractère de fin de ligne
'10'
>>> int(ligne1.rstrip()) #conversion de chaine vers entier
10
>>> ligne2 = g.readline()  #lecture de ligne 2
>>> ligne2
'11 12\n'
>>> ligne2.rstrip()  #nettoyage du caractère de fin de ligne
'11 12'
>>> ligne2.rstrip().split(' ')  #découpage de ligne 2 selon le caractère par défaut de séparation ' ', on peut écrire juste split() 
['11', '12']
>>> [int(c) for c in ligne2.rstrip().split()]  #conversion en entier
[11, 12]
>>> g.close()   #fermeture du fichier
\end{lstlisting}

\end{methode}

\begin{exercice}{}

\begin{itemize}
\tightlist
\item
  \textbf{Énoncé :}
\end{itemize}

Comme le disait le grand-père de Joseph Marchand, « Le plus beau voyage
est celui qu'on n'a pas encore fait ». New York était dans la tête de
Joseph depuis plusieurs années maintenant, et il a décidé aujourd'hui
d'acheter son billet d'avion. Quelques secondes avant de cliquer sur le
bouton « Acheter ! », son amie Haruhi lui envoie une liste de prix
qu'elle a trouvés sur Internet. Joseph est curieux de voir si ces
derniers sont moins chers que le billet qu'il s'apprêtait à acheter.

\begin{itemize}
\tightlist
\item
  \textbf{Entrée :}
\end{itemize}

Sur la première ligne le prix initial du billet de Joseph. Sur la
deuxième ligne un entier N , correspondant au nombre de billets envoyés
par Haruhi. La ligne suivante contient les N prix trouvés par Haruhi.

\begin{itemize}
\tightlist
\item
  \textbf{Sortie :}
\end{itemize}

Si Haruhi a trouvé au moins 3 prix strictement moins chers que celui de
Joseph, affichez « ARNAQUE ! » pour l'avertir. Sinon « Ok bon voyage,
bisous, n'oublie pas de m'envoyer des photos ! ».

\begin{itemize}
\item
  Exemple 1 :

  \begin{itemize}
  \item
    Entrée :

\begin{lstlisting}
570
4
495 1200 540 450
\end{lstlisting}
  \item
    Sortie :

\begin{lstlisting}
ARNAQUE !
\end{lstlisting}
  \end{itemize}
\item
  Exemple 2:

  \begin{itemize}
  \item
    Entrée :

\begin{lstlisting}
820
5
580 2000 970 1050 820
\end{lstlisting}
  \item
    Sortie :
  \end{itemize}

\begin{lstlisting}
Ok bon voyage, bisous, n'oublie pas de m'envoyer des photos !
\end{lstlisting}
\end{itemize}

\begin{enumerate}
\def\labelenumi{\arabic{enumi}.}
\item
  Récupérer les fichiers textes
  \passthrough{\lstinline!arnaque-exemple1.txt!} et
  \passthrough{\lstinline!arnaque-exemple2.txt!} avec les entrées des
  deux exemples.
\item
  En s'inspirant de la méthode de lecture de fichier texte présentée
  précédemment, compléter la fonction ci-dessous pour qu'elle résolve le
  problème. Les deux assertions doivent être vérifiées.
\item
  Se créer un compte sur le site
  \href{https://prologin.org/user/register?next=\%2F}{Prologin} et
  soumettre sa solution au juge en ligne sur
  \url{https://prologin.org/train/2019/qualification/arnaque_aerienne}.

  Avant de soumettre, remplacer dans la fonction l'instruction
  \passthrough{\lstinline!f = open(inputfile)!} par
  \passthrough{\lstinline!f = inputfile!} et le code du programme
  principal par :

\begin{lstlisting}[language=Python]
import sys
arnaque(sys.stdin)
\end{lstlisting}

  \passthrough{\lstinline!sys.stdin!} est
  \href{https://docs.python.org/fr/3/library/sys.html\#sys.stdin}{l'entrée
  standard} de l'interpréteur Python.
\item
  Sur la même plateforme, résoudre cet autre problème du même type :
  \url{https://prologin.org/train/2018/qualification/faites_place}
\end{enumerate}

\begin{lstlisting}[language=Python]
def arnaque(inputfile):
    """Résolution du problème 
    https://prologin.org/train/2019/qualification/arnaque_aerienne"""
    f = open(inputfile)  #ouverture du  fichier d'entrée
    # TO DO à compléter
    f.close() 
    c = 0
    # TO DO à compléter
    if c >= 3:
        print("ARNAQUE !")
        return 1
    else:
        print("Ok bon voyage, bisous, n'oublie pas de m'envoyer des photos !")
        return 0

# Tests unitaires        
assert arnaque('arnaque-exemple1.txt') == 1
assert arnaque('arnaque-exemple2.txt') == 0
print("Tests unitaires réussis !")
\end{lstlisting}

\end{exercice}

\begin{exercice}{}

Écrire une fonction \passthrough{\lstinline!maximum\_intervalle(t, n)!}
qui prend en paramètres un tableau d'entiers \passthrough{\lstinline!t!}
et un entier \passthrough{\lstinline!n!} supposé inférieur ou égal à la
longueur de \passthrough{\lstinline!t!} et qui retourne la somme
maximale sur tous les sous-tableaux de longueur
\passthrough{\lstinline!n!} inclus dans \passthrough{\lstinline!t!} .

\begin{lstlisting}[language=Python]
>>> maximum_intervalle([4,1,10,2,8,5], 3)
20
\end{lstlisting}

\end{exercice}

\end{document}
